\chapter{The Arithmetic--Geometric Mean}

\begin{center}\itshape
We have established that the arithmetic--geometric mean between $1$ and $\sqrt2$ is $\pi/\varpi$ to the eleventh decimal place; the
demonstration of this fact will surely open up an entirely new field of analysis.\\
\upshape --- Carl Friedrich Gauss, mathematical diary, 30 May 1799
\end{center}

Take two positive numbers, and replace them by their arithmetic mean and their geometric mean. Repeat. The two numbers approach
each other very fast: the number of correct digits doubles at every step. Their common limit, the \emph{arithmetic--geometric mean}, looks like
a curiosity. Gauss discovered at the age of twenty-one that it computes an integral, the one that gives the length of an ellipse, and through it the
number $\pi$. This chapter proves his theorem with a substitution found by Newman in 1985, computes the perimeter of an ellipse, proves
Legendre's relation between elliptic integrals, and combines everything into the algorithm of Brent and Salamin (1976), which computes $\pi$
with the number of correct digits doubling at every step. It closes the tome: the Gamma function, Wallis's product, Stirling, the lemniscate of
Chapter~5 and the fast multiplication of Chapter~7 all meet here.

\begin{think}
Start from $1$ and $\sqrt2$. After one step you have $1.2071$ and $1.1892$; after four, both are $1.19814023473559\ldots$ How many steps would you
expect to need for $1000$ digits? Gauss compared this number with $2\int_0^1\frac{\mathrm dt}{\sqrt{1-t^4}}=2.62205755\ldots$, the length of a half-lemniscate
(Chapter~5, Exercise~\ref{ex:lemniscate}). Compute $\pi$ divided by it.
\end{think}

\section{The arithmetic--geometric mean}

Let $a\ge b>0$. Define $a_0=a$, $b_0=b$ and
\[
\begin{gathered}a_{n+1}=\frac{a_n+b_n}2,\qq b_{n+1}=\sqrt{a_nb_n},\end{gathered}
\]
and $c_n=\sqrt{a_n^2-b_n^2}\ge0$.

\begin{proposition}\label{prop:agm}
\begin{enumerate}[label=(\alph*)]
\item $b\le b_n\le b_{n+1}\le a_{n+1}\le a_n\le a$ for every $n$, and $(a_n)$, $(b_n)$ have a common limit $M(a,b)$, the \emph{arithmetic--geometric mean} of $a$ and $b$.
\item $c_{n+1}=\frac{a_n-b_n}2$ and $c_{n+1}=\frac{c_n^2}{4a_{n+1}}$. Moreover $0\le a_{n+1}-b_{n+1}\le\frac{(a_n-b_n)^2}{8b}$.
\item $M(\lambda a,\lambda b)=\lambda M(a,b)$ for $\lambda>0$, and $M(a,b)=M\bigl(\frac{a+b}2,\sqrt{ab}\bigr)$.
\end{enumerate}
\end{proposition}
\begin{proof}
(a) The arithmetic--geometric inequality $\sqrt{a_nb_n}\le\frac{a_n+b_n}2$ gives $b_{n+1}\le a_{n+1}$, and since $b_n\le a_n$, both means lie between $b_n$ and $a_n$.
So $(a_n)$ decreases, $(b_n)$ increases, and both are bounded; they converge, to $A$ and $B$ say. Passing to the limit in $a_{n+1}=\frac{a_n+b_n}2$ gives $A=\frac{A+B}2$, so $A=B$.

(b) $c_{n+1}^2=\frac{(a_n+b_n)^2}4-a_nb_n=\frac{(a_n-b_n)^2}4$, and $\frac{a_n-b_n}2=\frac{a_n^2-b_n^2}{2(a_n+b_n)}=\frac{c_n^2}{4a_{n+1}}$. Finally
$a_{n+1}-b_{n+1}=\frac{(\sqrt{a_n}-\sqrt{b_n})^2}2=\frac{(a_n-b_n)^2}{2(\sqrt{a_n}+\sqrt{b_n})^2}$, and $(\sqrt{a_n}+\sqrt{b_n})^2\ge4b_n\ge4b$.

(c) The sequences built from $(\lambda a,\lambda b)$ are $(\lambda a_n,\lambda b_n)$. The sequences built from $(a_1,b_1)$ are those built from $(a,b)$, shifted by one.
\end{proof}

Part (b) is the reason for the speed. Once $a_n-b_n$ is small compared with $8b$, the next difference is of the order of its square: if $a_n-b_n\approx10^{-k}$,
then $a_{n+1}-b_{n+1}\approx10^{-2k}$. For $M(1,\sqrt2)$ the differences are $0.41$, $0.018$, $3.3\cdot10^{-5}$, $1.2\cdot10^{-10}$, $1.4\cdot10^{-21}$, $2\cdot10^{-43}$:
each step doubles the number of correct digits, and ten steps give a thousand digits. Before that, while $a_n$ is much larger than $b_n$, the
difference only halves at each step (Figure~\ref{fig:agm}).

\begin{figure}[ht]
\centering
\fig{ch12-agm}
\caption{The AGM iteration from $a_0=1$, $b_0=0.01$. The differences $a_n-b_n$ are $0.99$, $0.405$, $0.078$, $2.9\cdot10^{-3}$, $4.0\cdot10^{-6}$, $7.5\cdot10^{-12}$: slow
while the two numbers have different sizes, then quadratic. The limit is $M(1,0.01)=0.2621668872\ldots$}
\label{fig:agm}
\end{figure}

\begin{think}
Why does squaring the error double the number of correct digits? Newton's method for $\sqrt2$ has the same property; what do the two
algorithms have in common? (Look at $a_{n+1}-b_{n+1}$ as a function of $a_n-b_n$ near $0$.)
\end{think}

\section{Gauss's theorem}

For $a,b>0$ let
\[
I(a,b)=\int_0^{\pi/2}\frac{\mathrm d\theta}{\sqrt{a^2\cos^2\theta+b^2\sin^2\theta}}.
\]
The substitution $\theta\mapsto\frac\pi2-\theta$ shows $I(a,b)=I(b,a)$, $I(\lambda a,\lambda b)=\frac1\lambda I(a,b)$, and $I(a,a)=\frac\pi{2a}$.

\begin{theorem}[Gauss, 1799]\label{thm:gaussagm}
For $a,b>0$,
\[
I(a,b)=\frac\pi{2M(a,b)}.
\]
\end{theorem}
\begin{proof}
\emph{Step 1: an algebraic form.} Substitute $t=b\tan\theta$, a bijection from $\left[0,\frac\pi2\right[$ onto $[0,\infty[$, with $\mathrm dt=b(1+\tan^2\theta)\dd\theta$.
Then $a^2\cos^2\theta+b^2\sin^2\theta=\cos^2\theta\,(a^2+t^2)$ and $1+\tan^2\theta=\frac1{\cos^2\theta}=\frac{b^2+t^2}{b^2}$, so
\[
\begin{aligned}I(a,b)&=\int_0^\infty\frac{\mathrm dt}{\sqrt{(a^2+t^2)(b^2+t^2)}}\eqbr=\frac12\int_{-\infty}^{\infty}\frac{\mathrm dt}{\sqrt{(a^2+t^2)(b^2+t^2)}}.\end{aligned}
\]

\emph{Step 2: Newman's substitution.} Let $a_1=\frac{a+b}2$, $b_1=\sqrt{ab}$, and substitute $t=\frac12\bigl(x-\frac{ab}x\bigr)$, an increasing bijection from $]0,\infty[$
onto $\R$, with $\mathrm dt=\frac12\bigl(1+\frac{ab}{x^2}\bigr)\mathrm dx=\frac{x^2+ab}{2x^2}\dd x$. Expanding the squares,
\[
\begin{gathered}t^2+b_1^2=\frac{(x^2+ab)^2}{4x^2},\qq t^2+a_1^2=\frac{(x^2+a^2)(x^2+b^2)}{4x^2}.\end{gathered}
\]
(For the second: $t^2+a_1^2=\frac14\bigl(x^2-2ab+\frac{a^2b^2}{x^2}+a^2+2ab+b^2\bigr)$.) Write $D=\sqrt{(x^2+a^2)(x^2+b^2)}$. Then
$\sqrt{(a_1^2+t^2)(b_1^2+t^2)}=\frac{(x^2+ab)D}{4x^2}$, so
\[
\begin{aligned}\frac{\mathrm dt}{\sqrt{(a_1^2+t^2)(b_1^2+t^2)}}&=\frac{x^2+ab}{2x^2}\cdot\frac{4x^2}{(x^2+ab)D}\dd x\eqbr=\frac{2\dd x}D.\end{aligned}
\]
By Step~1, $I(a_1,b_1)=\frac12\int_0^\infty\frac{2\,\mathrm dx}{\sqrt{(x^2+a^2)(x^2+b^2)}}=I(a,b)$.

\emph{Step 3: the limit.} By induction $I(a,b)=I(a_n,b_n)$ for every $n$, with the sequences of Proposition~\ref{prop:agm} (after swapping $a$ and $b$ if
needed). Since $b_n^2\le a_n^2\cos^2\theta+b_n^2\sin^2\theta\le a_n^2$, we have $\frac\pi{2a_n}\le I(a_n,b_n)\le\frac\pi{2b_n}$, and both bounds tend to $\frac\pi{2M(a,b)}$.
\end{proof}

The proof is short because all the work is in Step~2, which is just the right change of variable. Gauss's own proof used a more complicated trigonometric substitution; the one above appears in a short paper of D.~J.~Newman (1985).

In terms of the modulus $k\in[0,1[$ and the complementary modulus $k'=\sqrt{1-k^2}$, the \emph{complete elliptic integrals} of the first and second kind are
\[
\begin{gathered}K(k)=\int_0^{\pi/2}\frac{\mathrm d\theta}{\sqrt{1-k^2\sin^2\theta}},\qq E(k)=\int_0^{\pi/2}\sqrt{1-k^2\sin^2\theta}\dd\theta.\end{gathered}
\]
Since $1-k^2\sin^2\theta=\cos^2\theta+k'^2\sin^2\theta$, we have $K(k)=I(1,k')=\frac\pi{2M(1,k')}$.

\begin{corollary}[Landen's transformation]\label{cor:landen}
For $0\le k<1$, let $k_1=\frac{1-k'}{1+k'}$. Then $K(k)=(1+k_1)K(k_1)$.
\end{corollary}
\begin{proof}
By Step~2 above and homogeneity, $K(k)=I(1,k')=I\bigl(\frac{1+k'}2,\sqrt{k'}\bigr)=\frac2{1+k'}I\bigl(1,\frac{2\sqrt{k'}}{1+k'}\bigr)$. The last argument is the complementary
modulus of $k_1$, since $1-\frac{4k'}{(1+k')^2}=\frac{(1-k')^2}{(1+k')^2}$. And $1+k_1=\frac2{1+k'}$.
\end{proof}

Since $k_1\approx\frac{k^2}4$ for small $k$, iterating Landen's transformation drives the modulus to $0$ quadratically, where $K(0)=\frac\pi2$: this is the AGM again,
seen through the modulus.

\begin{corollary}[the lemniscate]\label{cor:lemniscate}
Let $\varpi=2\int_0^1\frac{\mathrm dr}{\sqrt{1-r^4}}$, so that the lemniscate $(x^2+y^2)^2=x^2-y^2$ has length $2\varpi$ (Chapter~5). Then
\[
\begin{gathered}\frac\varpi\pi=\frac1{M(1,\sqrt2)},\qq\varpi=2.62205755429211981\ldots\end{gathered}
\]
\end{corollary}
\begin{proof}
Substitute $r=\cos\theta$: $\mathrm dr=-\sin\theta\dd\theta$ and $1-\cos^4\theta=\sin^2\theta\,(1+\cos^2\theta)$, so
$\int_0^1\frac{\mathrm dr}{\sqrt{1-r^4}}=\int_0^{\pi/2}\frac{\mathrm d\theta}{\sqrt{1+\cos^2\theta}}=\int_0^{\pi/2}\frac{\mathrm d\theta}{\sqrt{2\cos^2\theta+\sin^2\theta}}=I(\sqrt2,1)$. By Theorem~\ref{thm:gaussagm} this is
$\frac\pi{2M(1,\sqrt2)}$.
\end{proof}

\begin{figure}[ht]
\centering
\fig{ch12-lemniscate}
\caption{The lemniscate $(x^2+y^2)^2=x^2-y^2$, in polar coordinates $r^2=\cos2\theta$. Along it, $\mathrm ds=\frac{\mathrm dr}{\sqrt{1-r^4}}$, so the red arc from the centre to
$(1,0)$ has length $\int_0^1\frac{\mathrm dr}{\sqrt{1-r^4}}=\frac\varpi2$, and the whole curve has length $2\varpi$, as the circle of diameter $1$ has length $\pi$.}
\label{fig:lemniscate}
\end{figure}

\begin{beauty}
This is the identity of the epigraph. Gauss had computed both numbers to eleven decimals, and his diary entry records the moment he saw
they agreed, before he had a proof. He then developed the theory of the lemniscatic functions, the inverses of $\int\frac{\mathrm dr}{\sqrt{1-r^4}}$, which are
the first elliptic functions; the ``entirely new field of analysis'' he announced is the theory of elliptic functions and modular forms, the subject
of Tome~12. With Chapter~5, the corollary also computes $\Gamma(\frac14)$ by the AGM (Exercise~\ref{ex:gamma14}).
\end{beauty}

\section{The perimeter of an ellipse}

The ellipse with semi-axes $a\ge b>0$, parametrised by $(a\cos t,b\sin t)$, has speed $\sqrt{a^2\sin^2t+b^2\cos^2t}$, so its perimeter is $L(a,b)=4J(a,b)$, where
\[
J(a,b)=\int_0^{\pi/2}\sqrt{a^2\cos^2\theta+b^2\sin^2\theta}\dd\theta
\]
(substitute $\theta=\frac\pi2-t$). With $k$ the eccentricity, $b=ak'$ and $J(a,b)=aE(k)$. For a circle, $J(a,a)=\frac{\pi a}2$. In the rest of this section, $a,b>0$, we write
$R=a^2\cos^2\theta+b^2\sin^2\theta$, so that $I=\int_0^{\pi/2}R^{-1/2}$ and $J=\int_0^{\pi/2}R^{1/2}$, and the AGM sequences start from $a_0=a$, $b_0=b$.

\begin{lemma}\label{lem:jderiv}
\begin{enumerate}[label=(\alph*)]
\item $\displaystyle J(a,b)=a^2b^2\int_0^{\pi/2}R^{-3/2}\dd\theta$.
\item $I$ is of class $C^1$ on $]0,\infty[^2$, $a\partial_aI+b\partial_bI=-I$, and, for $a\ne b$,
\[
J(a,b)=a^2I(a,b)+a(a^2-b^2)\,\partial_aI(a,b).
\]
\end{enumerate}
\end{lemma}
\begin{proof}
(a) Write $c=\cos^2\theta$ and $s=\sin^2\theta$. The derivative of $\sin\theta\cos\theta\,R^{-1/2}$ is
$(c-s)R^{-1/2}-\sin\theta\cos\theta\cdot\frac12R^{-3/2}\cdot2(b^2-a^2)\sin\theta\cos\theta$, that is $\bigl((c-s)R-(b^2-a^2)cs\bigr)R^{-3/2}$. Using $c+s=1$, the
numerator is $a^2c^2-b^2s^2$, and so is $\frac{R^2-a^2b^2}{a^2-b^2}$ when $a\ne b$, because
$R^2-a^2b^2=a^4c^2+b^4s^2-a^2b^2(c^2+s^2)=(a^2-b^2)(a^2c^2-b^2s^2)$. So for $a\ne b$
\[
\frac{\mathrm d}{\mathrm d\theta}\Bigl(\frac{\sin\theta\cos\theta}{\sqrt R}\Bigr)=\frac{R^{1/2}-a^2b^2R^{-3/2}}{a^2-b^2}.
\]
The left side integrates to $0$ over $[0,\frac\pi2]$, which gives (a). For $a=b$, both sides of (a) equal $\frac{\pi a}2$.

(b) The integrand of $I$ is $C^1$ in $(a,b)$, with partial derivatives continuous on $[0,\frac\pi2]\times\left]0,\infty\right[^2$, so Theorem~\ref{thm:paramcompact} (applied in each variable)
gives $\partial_aI=-a\int\cos^2\theta\,R^{-3/2}$ and $\partial_bI=-b\int\sin^2\theta\,R^{-3/2}$, both continuous. Let $P=\int c\,R^{-3/2}$ and $Q=\int s\,R^{-3/2}$. Then
$a^2P+b^2Q=\int R\cdot R^{-3/2}=I$, which is the identity $a\partial_aI+b\partial_bI=-I$, and $P+Q=\int R^{-3/2}=\frac J{a^2b^2}$ by (a). Eliminating $Q$:
$(a^2-b^2)P=I-\frac J{a^2}$. Since $\partial_aI=-aP$, this is the formula.
\end{proof}

\begin{proposition}\label{prop:jstep}
For $a,b>0$, with $a_1=\frac{a+b}2$ and $b_1=\sqrt{ab}$,
\[
2J(a_1,b_1)-J(a,b)=ab\,I(a,b).
\]
\end{proposition}
\begin{proof}
Both sides are continuous in $(a,b)$, so it is enough to treat $a\ne b$. Write $I_1=\partial_1I(a_1,b_1)$ for the partial derivative of $I$ in its first variable at
$(a_1,b_1)$, and $I=I(a,b)=I(a_1,b_1)$. Differentiate the identity $I(a,b)=I\bigl(\frac{a+b}2,\sqrt{ab}\bigr)$ of Theorem~\ref{thm:gaussagm} in $a$, by the chain rule:
\[
\partial_aI(a,b)=\frac12I_1+\frac{b_1}{2a}\,\partial_2I(a_1,b_1).
\]
By Lemma~\ref{lem:jderiv}(b) at $(a_1,b_1)$, $b_1\partial_2I(a_1,b_1)=-I-a_1I_1$, so $\partial_aI(a,b)=\frac{a-a_1}{2a}I_1-\frac I{2a}=\frac{a-b}{4a}I_1-\frac I{2a}$. Put this
into Lemma~\ref{lem:jderiv}(b) at $(a,b)$:
\[
J(a,b)=a^2I-\frac{a^2-b^2}2I+\frac{(a^2-b^2)(a-b)}4I_1.
\]
Lemma~\ref{lem:jderiv}(b) at $(a_1,b_1)$, where $a_1^2-b_1^2=\frac{(a-b)^2}4$, gives $J(a_1,b_1)=a_1^2I+\frac{(a+b)(a-b)^2}8I_1$. So the last term above is
$2\bigl(J(a_1,b_1)-a_1^2I\bigr)$, and
\[
\begin{aligned}J(a,b)&=2J(a_1,b_1)+\Bigl(a^2-\frac{a^2-b^2}2-\frac{(a+b)^2}2\Bigr)I\\&=2J(a_1,b_1)-abI.\end{aligned}
\]
\end{proof}

\begin{theorem}[the perimeter by the AGM]\label{thm:perimeter}
For $a\ge b>0$, with $c_0^2=a^2-b^2$ and $c_n$ as in Proposition~\ref{prop:agm},
\[
J(a,b)=\Bigl(a^2-\sum_{n\ge0}2^{n-1}c_n^2\Bigr)I(a,b),
\]
so the ellipse with semi-axes $a\ge b$ has perimeter
\[
L(a,b)=\frac{2\pi}{M(a,b)}\Bigl(a^2-\sum_{n\ge0}2^{n-1}c_n^2\Bigr).
\]
In terms of the modulus, with the AGM started from $1$ and $k'$, so that $c_0=k$,
\[
E(k)=\Bigl(1-\sum_{n\ge0}2^{n-1}c_n^2\Bigr)K(k).
\]
\end{theorem}
\begin{proof}
Let $I=I(a,b)=I(a_n,b_n)$ for all $n$, and $d_n=J(a_n,b_n)-a_n^2I$. Proposition~\ref{prop:jstep} at $(a_n,b_n)$ reads $J(a_n,b_n)=2J(a_{n+1},b_{n+1})-a_nb_nI$, that is
\[
\begin{aligned}d_n&=2d_{n+1}+\bigl(2a_{n+1}^2-a_nb_n-a_n^2\bigr)I\eqbr=2d_{n+1}-\frac{c_n^2}2I,\end{aligned}
\]
because $2a_{n+1}^2=\frac{a_n^2+2a_nb_n+b_n^2}2$. By induction, $d_0=2^Nd_N-I\sum_{n<N}2^{n-1}c_n^2$. Now $\frac\pi2b_n\le J(a_n,b_n)\le\frac\pi2a_n$ and $\frac\pi2a_n\le a_n^2I\le\frac{\pi a_n^2}{2b_n}$, so $|d_n|\le\frac{\pi(a_n^2-b_n^2)}{2b_n}\le\frac{\pi c_n^2}{2b}$. By Proposition~\ref{prop:agm}(b), $c_{n+1}\le\frac{c_n^2}{4M}$ with $M=M(a,b)$; since $c_n\to0$, there is $N_0$ with $\frac{c_{N_0}}{4M}\le\frac12$, and then
$c_{N_0+j}\le4M\cdot2^{-2^j}$. So $2^Nc_N^2\to0$, $2^Nd_N\to0$, and the series converges: $d_0=-I\sum_{n\ge0}2^{n-1}c_n^2$, which is the formula. The perimeter is
$4J$ and $I=\frac\pi{2M}$. For the last statement take $a=1$, $b=k'$.
\end{proof}

For the ellipse with semi-axes $2$ and $1$, the AGM gives $L=9.68844822054767619842\ldots$ after five steps, with every digit correct. For the Earth's orbit,
with eccentricity $0.0167$, the perimeter is $2\pi a\times0.99993027\ldots$: an ellipse that close to a circle is a circle for most purposes.

\begin{figure}[ht]
\centering
\fig{ch12-perimeter}
\caption{The perimeter of the ellipse with semi-axes $1$ and $\sqrt{1-e^2}$, divided by $2\pi$, against the eccentricity $e$ (blue): it falls from $1$ (the circle) to
$\frac2\pi$ (the segment $[-1,1]$, travelled twice). Dashed: Euler's bounds $\frac{\pi(a+b)}{2\pi}$ (green) and $\frac{\pi\sqrt{2(a^2+b^2)}}{2\pi}$ (orange) of Exercise~\ref{ex:eulerbounds}.}
\label{fig:perimeter}
\end{figure}

\begin{think}
The area of an ellipse is $\pi ab$. Is there a formula of the same kind for its perimeter, built from $a$, $b$, square roots, $\pi$, logarithms and the
like? What would such a formula give when $b\to0$, and when $b=a$?
\end{think}

\begin{remark}
There is none: $E(k)$ is not an elementary function of $k$. This follows from Liouville's theory of integration in finite terms, which decides when an
integral of an algebraic function is elementary; see the survey of Marchisotto and Zakeri in the references. So Theorem~\ref{thm:perimeter}, a limit of
algebraic operations converging quadratically, is as close to a closed form as one can get. Ramanujan found remarkable approximations, such as
$L\approx\pi\bigl(3(a+b)-\sqrt{(3a+b)(a+3b)}\bigr)$, whose relative error is $3\cdot10^{-6}$ for $a=2b$ and at most $0.5\%$ for all ellipses; a numerical calculation tests them
against the AGM.
\end{remark}

\section{Legendre's relation and the computation of $\pi$}

\begin{proposition}\label{prop:kederiv}
On $]0,1[$, $K$ and $E$ are differentiable, with
\[
\begin{gathered}\frac{\mathrm dE}{\mathrm dk}=\frac{E-K}k,\qq\frac{\mathrm dK}{\mathrm dk}=\frac{E-k'^2K}{k\,k'^2}.\end{gathered}
\]
\end{proposition}
\begin{proof}
Let $\Delta=\sqrt{1-k^2\sin^2\theta}$. Differentiating under the integral sign (Theorem~\ref{thm:paramcompact}, on $[0,\frac\pi2]\times[0,k_1]$ for any $k_1<1$),
$\frac{\mathrm dE}{\mathrm dk}=-\int\frac{k\sin^2\theta}\Delta$ and $\frac{\mathrm dK}{\mathrm dk}=\int\frac{k\sin^2\theta}{\Delta^3}$. Since $k^2\sin^2\theta=1-\Delta^2$, the first is
$\frac1k\int\frac{\Delta^2-1}\Delta=\frac{E-K}k$, and the second is $\frac1k\bigl(\int\Delta^{-3}-K\bigr)$. Lemma~\ref{lem:jderiv}(a) with $a=1$, $b=k'$ (so that $R=\Delta^2$ and
$J=E$) gives $\int\Delta^{-3}=\frac E{k'^2}$.
\end{proof}

\begin{theorem}[Legendre's relation, 1811]\label{thm:legendre}
For $0<k<1$, with $K'=K(k')$ and $E'=E(k')$,
\[
EK'+E'K-KK'=\frac\pi2.
\]
\end{theorem}
\begin{proof}
Let $\Lambda(k)=EK'+E'K-KK'$, a function of $k$ alone since $k'=\sqrt{1-k^2}$. By Proposition~\ref{prop:kederiv} and the chain rule with $\frac{\mathrm dk'}{\mathrm dk}=-\frac k{k'}$,
\[
\begin{gathered}\frac{\mathrm dE'}{\mathrm dk}=-\frac{k(E'-K')}{k'^2},\qq\frac{\mathrm dK'}{\mathrm dk}=-\frac{E'-k^2K'}{k\,k'^2}.\end{gathered}
\]
Multiply $\Lambda'=E_kK'+EK'_k+E'_kK+E'K_k-K_kK'-KK'_k$ by $kk'^2$ and substitute. The terms in $EE'$ are $-EE'+E'E=0$. The terms in $EK'$ are
$k'^2+k^2-1=0$ times $EK'$; those in $E'K$ are $(-k^2-k'^2+1)E'K=0$; and those in $KK'$ are $(-k'^2+k^2+k'^2-k^2)KK'=0$. So $\Lambda$ is constant.

To find the constant, let $k\to0^+$. First, $K-E=\int\frac{k^2\sin^2\theta}\Delta\le\frac{\pi k^2}{2k'}$. Second, $K'=\int_0^{\pi/2}\frac{\mathrm d\theta}{\sqrt{\cos^2\theta+k^2\sin^2\theta}}=O\bigl(\ln\frac1k\bigr)$:
with $u=\frac\pi2-\theta$, on $u\le\frac\pi4$ the integrand is at most $\bigl(\frac{4u^2}{\pi^2}+\frac{k^2}2\bigr)^{-1/2}$, whose integral is $\frac\pi2\operatorname{arsinh}\frac1{\sqrt2\,k}$, and on $u\ge\frac\pi4$
it is at most $\sqrt2$. Third, $E'=\int\sqrt{\cos^2\theta+k^2\sin^2\theta}\to\int_0^{\pi/2}\cos\theta\dd\theta=1$, since the integrands differ by at most $k$. Hence
$\Lambda=E'K-(K-E)K'\to1\cdot\frac\pi2-0$.
\end{proof}

\begin{theorem}[Brent--Salamin, 1976]\label{thm:brentsalamin}
Run the AGM from $a_0=1$, $b_0=\frac1{\sqrt2}$, and let $c_n=\frac{a_{n-1}-b_{n-1}}2$ for $n\ge1$. Then
\[
\pi=\frac{4M(1,\frac1{\sqrt2})^2}{1-\sum_{n\ge1}2^{n+1}c_n^2},
\]
and the approximations $\pi_n=\frac{4a_{n+1}^2}{1-\sum_{j=1}^{n}2^{j+1}c_j^2}$ are correct to $1$, $4$, $9$, $20$, $42$, $85$, $173$, $347$ digits for $n=1,\dots,8$.
\end{theorem}
\begin{proof}
At $k=k'=\frac1{\sqrt2}$, Legendre's relation reads $2EK-K^2=\frac\pi2$. By Theorems~\ref{thm:gaussagm} and~\ref{thm:perimeter}, with $M=M(1,\frac1{\sqrt2})$, $K=\frac\pi{2M}$ and
$E=\bigl(1-\sum_{n\ge0}2^{n-1}c_n^2\bigr)K$, where $c_0^2=\frac12$ (and $c_n=\frac{a_{n-1}-b_{n-1}}2$ for $n\ge1$ by Proposition~\ref{prop:agm}(b)). So
\[
\frac\pi2=K^2\Bigl(1-\sum_{n\ge0}2^nc_n^2\Bigr)=\frac{\pi^2}{4M^2}\Bigl(\frac12-\sum_{n\ge1}2^nc_n^2\Bigr),
\]
which rearranges to the formula. The digit counts are computed numerically.
\end{proof}

\begin{remark}
Each step of the algorithm costs a few multiplications and one square root of numbers with $N$ digits. With the fast multiplication of Chapter~7 and
Newton's method for square roots, both cost $O(N\log N\log\log N)$ operations, and about $\log_2N$ steps are needed, so $N$ digits of $\pi$ cost
$O(N(\log N)^2\log\log N)$: far less than the $O(N^2)$ of an arctangent series such as Machin's $\frac\pi4=4\arctan\frac15-\arctan\frac1{239}$ with schoolbook arithmetic. Salamin and Brent
found the algorithm independently in 1976; every ingredient had been known to Gauss and Legendre. Since then, most records have used instead the series
of the Chudnovsky brothers (1988), $\frac1\pi=12\sum_{n\ge0}\frac{(-1)^n(6n)!\,(13591409+545140134n)}{(3n)!\,(n!)^3\,640320^{3n+3/2}}$, which gives $14$ digits per term and is
evaluated by binary splitting with fast multiplication. Its proof uses modular forms, the subject of Tome~12. The three methods can be compared numerically.
\end{remark}

\section{Exercises}

\begin{exercise}[$\star$]\label{ex:agmsteps}
Let $a=1$, $b=\frac1{\sqrt2}$. Using Proposition~\ref{prop:agm}(b), show that $a_n-b_n\le8b\,\bigl(\frac{a-b}{8b}\bigr)^{2^n}$, and deduce that $10$ steps give $M(1,\frac1{\sqrt2})$ to
more than $1000$ digits.
\end{exercise}

\begin{exercise}[$\star$]\label{ex:eulerbounds}
Show that the perimeter of an ellipse with semi-axes $a,b$ satisfies $\pi(a+b)\le L(a,b)\le\pi\sqrt{2(a^2+b^2)}$, with equality only for a circle.
\end{exercise}

\begin{exercise}[$\star\star$]\label{ex:pendulum}
A pendulum of length $\ell$ released at rest from the angle $\theta_0\in\left]0,\pi\right[$ satisfies $\ddot\theta=-\frac g\ell\sin\theta$. Show that its period is
\[
T=4\sqrt{\frac\ell g}\,K\Bigl(\sin\frac{\theta_0}2\Bigr)=\frac{2\pi\sqrt{\ell/g}}{M\bigl(1,\cos\frac{\theta_0}2\bigr)}.
\]
By how much is the period at $\theta_0=\frac\pi2$ longer than the small-oscillation period $2\pi\sqrt{\ell/g}$?
\end{exercise}

\begin{exercise}[$\star\star$]\label{ex:gamma14}
Using Exercise~\ref{ex:lemniscate} of Chapter~5 and Corollary~\ref{cor:lemniscate}, show that
\[
\Gamma\Bigl(\frac14\Bigr)=\frac{(2\pi)^{3/4}}{\sqrt{M(1,\sqrt2)}},
\]
and compute $\Gamma(\frac14)$ to ten digits with five steps of the AGM.
\end{exercise}

\begin{exercise}[$\star\star\star$]\label{ex:agmlog}
Show that, as $x\to0^+$,
\[
\frac\pi{2M(1,x)}=\ln\frac4x+o(1).
\]
So the AGM computes logarithms, with the same quadratic speed.
\end{exercise}

\section{Hints}

\begin{hint}{ex:agmsteps}
Put $e_n=\frac{a_n-b_n}{8b}$; then $e_{n+1}\le e_n^2$.
\end{hint}
\begin{hint}{ex:eulerbounds}
For the lower bound, compare the integrand $\sqrt{a^2\cos^2\theta+b^2\sin^2\theta}$ with the weighted mean $a\cos^2\theta+b\sin^2\theta$; for the upper bound use the Cauchy--Schwarz inequality.
\end{hint}
\begin{hint}{ex:pendulum}
Conservation of energy gives $\dot\theta^2$ as a function of $\theta$. Then substitute $\sin\frac\theta2=\sin\frac{\theta_0}2\sin\varphi$.
\end{hint}
\begin{hint}{ex:gamma14}
Both results compute the same integral $\int_0^1\frac{\mathrm dr}{\sqrt{1-r^4}}$.
\end{hint}
\begin{hint}{ex:agmlog}
Use the algebraic form $I(1,x)=\int_0^\infty\frac{\mathrm dt}{\sqrt{(1+t^2)(x^2+t^2)}}$ and split the integral at $t=1$.
\end{hint}

\section{Solutions}

\begin{solution}{ex:agmsteps}
\idea Quadratic convergence means $e_{n+1}\le e_n^2$ for a suitably normalised error, and then $e_n\le e_0^{2^n}$: the exponent doubles.
\writeup Proposition~\ref{prop:agm}(b) gives $a_{n+1}-b_{n+1}\le\frac{(a_n-b_n)^2}{8b}$, that is $e_{n+1}\le e_n^2$, so $e_n\le e_0^{2^n}$ by induction. Here
$e_0=\frac{1-0.7071}{8\cdot0.7071}\approx0.0518$, so $\log_{10}e_0\approx-1.286$ and $a_n-b_n\le5.66\cdot10^{-1.286\cdot2^n}$. Since $M$ lies between $b_n$ and $a_n$, the error of $a_n$
is below $a_n-b_n$. For $n=10$, $1.286\cdot1024>1300$. (The bound is pessimistic: the actual difference after $10$ steps is about $10^{-1396}$.)
\end{solution}

\begin{solution}{ex:eulerbounds}
\idea The perimeter is an average of $\sqrt{a^2\cos^2+b^2\sin^2}$; comparing a quadratic mean with an arithmetic mean gives one bound, and Cauchy--Schwarz the other.
\writeup With $c=\cos^2\theta$, $s=\sin^2\theta$, $c+s=1$: $(ac+bs)^2\le a^2c+b^2s$ by convexity of $x\mapsto x^2$ (with weights $c$, $s$), so
$\sqrt{a^2c+b^2s}\ge ac+bs$, and integrating over $[0,\frac\pi2]$, where $\int c=\int s=\frac\pi4$, gives $J\ge\frac{\pi(a+b)}4$. By Cauchy--Schwarz,
$J\le\sqrt{\frac\pi2}\sqrt{\int(a^2c+b^2s)}=\sqrt{\frac\pi2\cdot\frac{\pi(a^2+b^2)}4}$. Multiply by $4$. Equality in the first needs $ac+bs$ to be the quadratic mean for all
$\theta$, which forces $a=b$; in the second, the integrand must be constant, which again forces $a=b$. Figure~\ref{fig:perimeter} shows both bounds; the lower one is
much better for moderate eccentricity.
\end{solution}

\begin{solution}{ex:pendulum}
\idea Energy conservation turns the period into an integral $\int\frac{\mathrm d\theta}{\dot\theta}$, and a trigonometric substitution puts it in the form $K(k)$; Gauss's theorem
then gives the AGM.
\writeup Multiplying the equation by $\dot\theta$ and integrating, $\dot\theta^2=\frac{2g}\ell(\cos\theta-\cos\theta_0)=\frac{4g}\ell\bigl(k^2-\sin^2\frac\theta2\bigr)$ with $k=\sin\frac{\theta_0}2$.
By symmetry, the period is four times the time from $\theta=0$ to $\theta_0$, so $T=4\int_0^{\theta_0}\frac{\mathrm d\theta}{\dot\theta}=2\sqrt{\frac\ell g}\int_0^{\theta_0}\frac{\mathrm d\theta}{\sqrt{k^2-\sin^2(\theta/2)}}$.
Substitute $\sin\frac\theta2=k\sin\varphi$, $\varphi\in[0,\frac\pi2]$: then $\frac12\cos\frac\theta2\dd\theta=k\cos\varphi\dd\varphi$, $\sqrt{k^2-\sin^2\frac\theta2}=k\cos\varphi$ and
$\cos\frac\theta2=\sqrt{1-k^2\sin^2\varphi}$, so
\[
T=4\sqrt{\frac\ell g}\int_0^{\pi/2}\frac{\mathrm d\varphi}{\sqrt{1-k^2\sin^2\varphi}}=4\sqrt{\frac\ell g}K(k).
\]
With $k'=\cos\frac{\theta_0}2$ and $K(k)=\frac\pi{2M(1,k')}$ this is the formula. At $\theta_0=\frac\pi2$, $M(1,\cos\frac\pi4)=M(1,0.7071)=0.8472\ldots$, so the period is $18.0\%$ longer.
(For $10^\circ$ it is $0.19\%$ longer; for $179^\circ$, almost four times as long.)
\end{solution}

\begin{solution}{ex:gamma14}
\idea Two different evaluations of the same integral: one by the Beta function (Chapter~5), one by the AGM (this chapter).
\writeup Chapter~5 gives $\int_0^1\frac{\mathrm dr}{\sqrt{1-r^4}}=\frac{\Gamma(1/4)^2}{4\sqrt{2\pi}}$, and Corollary~\ref{cor:lemniscate} gives $\frac\pi{2M(1,\sqrt2)}$. Equating,
$\Gamma(\frac14)^2=\frac{2\pi\sqrt{2\pi}}{M(1,\sqrt2)}=\frac{(2\pi)^{3/2}}{M(1,\sqrt2)}$. After five steps, $a_5=1.19814023473559220744$, already equal to $M(1,\sqrt2)$ to all twenty-one digits shown, and the formula gives
$\Gamma(\frac14)=3.62560990822\ldots$, which agrees with the value computed from Euler's integral.
\end{solution}

\begin{solution}{ex:agmlog}
\idea For small $x$, $\frac1{\sqrt{x^2+t^2}}$ is close to $\frac1t$ except near $0$, where it is cut off at $t\approx x$: this produces $\ln\frac1x$, and the constants come from
exact antiderivatives.
\writeup By Theorem~\ref{thm:gaussagm} and Step~1 of its proof, $\frac\pi{2M(1,x)}=I(1,x)=\int_0^\infty\frac{\mathrm dt}{\sqrt{(1+t^2)(x^2+t^2)}}$. Split at $t=1$.

On $[1,\infty[$, the integrand converges to $\frac1{t\sqrt{1+t^2}}$ as $x\to0$, uniformly, and is bounded by it, so by dominated convergence (Corollary~\ref{cor:dct}) that part tends to
$\int_1^\infty\frac{\mathrm dt}{t\sqrt{1+t^2}}=\operatorname{arsinh}1=\ln(1+\sqrt2)$, using the antiderivative $-\operatorname{arsinh}\frac1t$.

On $[0,1]$, write the integrand as $\frac1{\sqrt{x^2+t^2}}+\frac{g(t)}{\sqrt{x^2+t^2}}$ with $g(t)=\frac1{\sqrt{1+t^2}}-1$. First, $\int_0^1\frac{\mathrm dt}{\sqrt{x^2+t^2}}=\operatorname{arsinh}\frac1x
=\ln\frac{1+\sqrt{1+x^2}}x=\ln\frac2x+o(1)$. Second, $|g(t)|\le\frac{t^2}2$, so $\bigl|\frac{g(t)}{\sqrt{x^2+t^2}}\bigr|\le\frac t2$, and by Corollary~\ref{cor:dct} (the convergence is uniform, the difference being at most $\frac x2$)
$\int_0^1\frac{g(t)}{\sqrt{x^2+t^2}}\dd t\to\int_0^1\frac{g(t)}t\dd t$. This integral equals $\bigl[-\operatorname{arsinh}\frac1t-\ln t\bigr]_0^1=\bigl[-\ln\bigl(1+\sqrt{1+t^2}\bigr)\bigr]_0^1=\ln2-\ln(1+\sqrt2)$.

Adding the three limits, $I(1,x)=\ln\frac2x+\ln2-\ln(1+\sqrt2)+\ln(1+\sqrt2)+o(1)=\ln\frac4x+o(1)$. Numerically, one finds the error $o(1)$ to be $6.8\cdot10^{-3}$ at $x=0.1$, $1.25\cdot10^{-4}$
at $x=0.01$ and $3.6\cdot10^{-12}$ at $x=10^{-6}$, consistent with the known bound $O(x^2\ln\frac1x)$ (Borwein and Borwein, Chapter~7). Taking $x=2^{-m}$ gives $(m+2)\ln2\approx\frac\pi{2M(1,2^{-m})}$, with an error of order $m4^{-m}$: with $\pi$ from Brent--Salamin, this computes $\ln2$ to about $0.6m$ digits.
\end{solution}

\section*{Sources and further reading}

\begin{itemize}
\item J.~M.~Borwein and P.~B.~Borwein, \emph{Pi and the AGM}, Wiley, 1987. Chapter~1 (the AGM, Gauss's theorem, Legendre's relation and the
algorithm for $\pi$) follows the same route as this chapter; Chapter~7 gives the logarithm by the AGM with error bounds.
\item D.~A.~Cox, \emph{The arithmetic-geometric mean of Gauss}, L'Enseignement Math\'ematique 30 (1984), 275--330. The history, Gauss's diary,
and the complex AGM, whose many limits are explained by modular forms.
\item D.~J.~Newman, \emph{A simplified version of the fast algorithms of Brent and Salamin}, Mathematics of Computation 44 (1985), 207--210.
\item E.~Salamin, \emph{Computation of $\pi$ using arithmetic-geometric mean}, Mathematics of Computation 30 (1976), 565--570; and R.~P.~Brent,
\emph{Fast multiple-precision evaluation of elementary functions}, Journal of the ACM 23 (1976), 242--251.
\item E.~A.~Marchisotto and G.-A.~Zakeri, \emph{An invitation to integration in finite terms}, College Mathematics Journal 25 (1994), 295--308.
\item D.~V.~Chudnovsky and G.~V.~Chudnovsky, \emph{Approximations and complex multiplication according to Ramanujan}, in \emph{Ramanujan
Revisited}, Academic Press, 1988, 375--472.
\end{itemize}
